\section{Part One: Two Layers of Vendor, and Why the Distinction Matters}

Before anything else can be reasoned about clearly, one terminological confusion needs resolving, because it silently collapses two different claims into one word.

\textbf{Model vendors} --- Anthropic, OpenAI, Google, DeepSeek, and similar organizations --- build large language models and their native harnesses: the decoding stack, safety training, tool-use format, default scaffolding (reasoning traces, and so on).

\textbf{Tooling vendors} build products around one or more models, wrapping them in an application-specific harness --- retrieval, verification logic, control flow, an action layer. A contract-review product or a customer-service agent is a tooling-vendor product in the large majority of cases.

The distinction is useful, but the boundary is not as clean in practice as a two-tier description suggests, and it's worth being precise about the actual range: some tooling vendors call a model vendor's API and do essentially nothing else; others self-host an open-weight model; others fine-tune or distill a model for their specific task; and a smaller number train a specialist model largely from scratch, sometimes combining several of these approaches within one product. \textbf{The more defensible general statement is that many AI products depend substantially on a third-party foundation model, although the form and depth of that dependency varies considerably} --- and the appropriate evaluation method shifts accordingly. A pure API wrapper's quality is dominated by model choice and harness design; a product built on a heavily fine-tuned or distilled model needs to be evaluated on that fine-tuning's own merits, separately from the base model's.

The distinction still matters for deciding \textbf{who is responsible for which failure and which claims need which kind of check.} A claim about raw model capability, safety testing, or training data is a \textbf{model-vendor claim}, best tested by independent benchmarking or red-teaming. A claim about how a specific product behaves --- "we verify every citation," "the agent runs in an isolated sandbox" --- is a \textbf{tooling-vendor claim}, best tested by a direct teardown of that product's actual architecture.

\textbf{One correction to how v1 stated the relationship between the two layers:} v1 claimed a tooling vendor "inherits every limitation" of the underlying model and "cannot patch" its failure modes. That overstates the case. What's true is narrower: a tooling vendor cannot change the \emph{unmodified base model's weights} --- the raw model, called with no scaffolding, has whatever failure profile it has. But a tooling vendor's system-level choices --- retrieval, constrained decoding, deterministic tool calls, targeted fine-tuning, routing between models, restricting what actions the system is allowed to take --- can eliminate or substantially reduce the \emph{deployed system's} failure rate on a defined task, even without changing the base model at all. The useful distinction is between \textbf{limitations of the unchanged base weights} (a model-vendor-layer property) and \textbf{the failure rate of the deployed system as actually built} (a tooling-vendor-layer property, which can differ enormously from the base model's raw behavior). Conflating the two erases exactly the leverage a tooling vendor has to offer.
